\chapter{Physical Examination by a Doctor}
\index{Physical Examination by a Doctor}

\medskip
\noindent\textit{This chapter is for general education. A physical examination is one part of an assessment and does not replace discussion of symptoms and concerns.}

\section*{What it is}
A physical examination is a structured assessment of the body by a trained healthcare professional. It may help identify the cause and severity of a problem, establish a baseline, guide tests, and monitor treatment.

\section*{What may happen}
With your permission, the clinician may observe how you look and move, measure vital signs, feel or press areas of the body, listen to the heart and lungs, examine the abdomen, assess nerves and movement, or inspect the skin, mouth, ears, and eyes. The examination should be explained and adapted to your symptoms, age, communication needs, disability, and cultural or privacy concerns.

\section*{Preparing and participating}
Bring a list of medicines, allergies, relevant records, and questions. Tell the clinician about pain, pregnancy possibility, past trauma, or anything that may make an examination difficult. You can ask for a chaperone, a support person where permitted, an interpreter, a different clinician, or a pause. You can decline or stop an examination, although this may limit what can be assessed.

\section*{After the examination}
The clinician may explain likely causes, recommend tests, give treatment, or arrange follow-up. Ask what the findings mean, when results will be available, and which symptoms should prompt urgent help. An examination cannot rule out every condition; further testing may be needed.

\section*{When to Seek Medical Care}
Seek urgent care for severe or rapidly worsening symptoms, collapse, severe breathing difficulty, chest pain, new weakness or confusion, or uncontrolled bleeding. Contact the clinic if an examination causes persistent severe pain, bleeding, or another concerning change.

\section*{Sources and further reading}
\begin{itemize}
\item MedlinePlus. \textit{Medical tests}. \url{https://medlineplus.gov/medicaltests.html}
\item World Health Organization. \textit{Quality health services}. \url{https://www.who.int/health-topics/quality-health-services}
\end{itemize}
